% Eigenständige Lesefassung. Nummerierte Aussagen und Beweistabellen stehen
% in Band 37; die Nullprodukt-Rekonstruktion besitzt eine eigene Lesefassung.
\chapter*{Endliche Halbgruppen erkennen}
\hypertarget{finite-semigroups.reading}{}
\addcontentsline{toc}{chapter}{Endliche Halbgruppen erkennen}

Eine endliche Multiplikationstafel lässt sich vollständig hinschreiben.
Damit ist aber noch nicht erklärt, welche Struktur sie besitzt oder ob
eine anders beschriftete Tafel dieselbe Struktur beschreibt. In diesem
Text verbinden wir drei Fragen: Was erzwingt die Endlichkeit für die
Potenzen eines Elements? Wie erkennt man dieselbe Halbgruppe unter
anderen Namen? Und welche zusätzliche Schwierigkeit entsteht, wenn
statt der Elemente nur die Multiplikation ihrer Teilmengen gegeben ist?

Eine \emph{Halbgruppe} ist eine Menge \(S\) mit einer Operation
\(S\times S\to S\), für die
\[
  (xy)z=x(yz)\qquad(x,y,z\in S)
\]
gilt. Die Ergebnisse liegen wieder in \(S\); wir können deshalb weiter
multiplizieren. Die Assoziativität erlaubt, bei längeren Produkten die
Klammern wegzulassen. Ein neutrales Element, ein Nullelement oder die
Vertauschbarkeit der Faktoren werden dabei nicht vorausgesetzt.
\emph{Endlich} bedeutet, dass \(S\) nur endlich viele Elemente besitzt.
Diese beiden Forderungen bilden die Definition
\FormulaRefAuto{FiniteSemigroupDef}[def] in Band~37.
Wenn wir ein Element oder ein nichtleeres Ideal wählen, verlangen wir
zusätzlich \(S\ne\varnothing\). Ein leerer Träger besitzt nur die
leere binäre Operation; seine gesonderte Behandlung ist unmittelbar.

\medskip\noindent\textbf{Welche Art von Antwort suchen wir?}
Bei gegebenen Grundtafeln erhalten wir einen vollständigen
Isomorphietest. Bei gegebenen Potenzhalbgruppen zeigen wir zunächst,
welche Informationen sicher zurückgewonnen werden können und wo dieser
Test eine weitere Begründung braucht. Für eine festgelegte
Grundelementzahl lässt sich die Rekonstruktionsfrage anschließend durch
eine endliche Suche entscheiden. Die Begrenzung auf diese eine Zahl
gehört zum Ergebnis der Suche dazu.

\section*{Ein Dreierbeispiel ohne Null}
\hypertarget{finite-semigroups.reading.example}{}

\begin{FiniteSemigroupsReadingExample}{Eine Potenzfolge mit Anlauf und Kreis}
Seien \(a,b,c\) verschieden und \(S=\{a,b,c\}\). Wir verwenden die Tafel
\[
  \begin{array}{c|ccc}
    \cdot&a&b&c\\\hline
    a&b&c&b\\
    b&c&b&c\\
    c&b&c&b
  \end{array}.
\]
Aus ihr folgen
\[
  a^1=a,\quad a^2=b,\quad a^3=c,\quad a^4=b,\quad
  a^5=c,\quad a^6=b,\ldots
\]
Nach dem ersten Wert \(a\) wechseln sich \(b\) und \(c\) ab.
Keines der drei Elemente ist eine Null: Schon \(a^2\ne a\),
\(bc\ne b\) und \(c^2\ne c\) schließen dies jeweils aus.
Auch ein neutrales Element für ganz \(S\) gibt es nicht:
Jedes mögliche Element scheitert bereits an der Multiplikation mit \(a\).
\end{FiniteSemigroupsReadingExample}

Bevor wir das Beispiel benutzen, beweisen wir seine Assoziativität.
Schreibe \(x_1=a,x_2=b,x_3=c\) und definiere für jeden positiven
Exponenten \(r\)
\[
  \rho(r)=
  \begin{cases}
    1,&r=1,\\
    2,&r\geq2\text{ und }r\text{ gerade},\\
    3,&r\geq3\text{ und }r\text{ ungerade}.
  \end{cases}
\]
Die Tafel besagt genau \(x_rx_s=x_{\rho(r+s)}\) für
\(r,s\in\{1,2,3\}\). Das Ersetzen einer Zahl \(t\) durch \(\rho(t)\)
erhält ihre Parität. Bei Summen mit einem weiteren positiven Exponenten
liegen außerdem beide Summen mindestens bei \(2\). Deshalb gilt
\[
  \rho\bigl(\rho(r+s)+t\bigr)=\rho(r+s+t)
    =\rho\bigl(r+\rho(s+t)\bigr).
\]
Beide Klammerungen von \(x_rx_sx_t\) liefern also denselben Wert.
Das beweist die Assoziativität der ganzen Tafel. Wir haben nicht nur
einzelne Produkte geprüft, sondern eine gemeinsame Regel für alle
\(3^3=27\) Tripel angegeben.

\section*{Warum sich Potenzen schließlich wiederholen}
\hypertarget{finite-semigroups.reading.powers}{}

Für ein Element \(a\in S\) sind die \emph{positiven Potenzen} rekursiv
definiert durch \(a^1=a\) und \(a^{r+1}=a^ra\).
Eine Potenz \(a^0\) brauchen wir nicht; ohne neutrales Element hätte sie
keinen vorgegebenen Wert. Aus der Assoziativität folgt das Additionsgesetz
\[
  a^{r+s}=a^ra^s\qquad(r,s\geq1).
\]
Es wird in Band~28 hergeleitet und ist die Grundlage für das Folgende.

Hat \(S\) genau \(n\) Elemente, können die \(n+1\) Werte
\(a^1,\ldots,a^{n+1}\) nicht alle verschieden sein. Es gibt also
positive Exponenten \(i<j\) mit \(a^i=a^j\). Das ist die Anwendung des
Schubfachprinzips in
\FormulaRefAuto{FiniteSemigroupPowerRepetitionExists}[thm].
Der endliche Wertebereich erzwingt eine Wiederholung, obwohl die
Exponenten unbegrenzt weiterlaufen.

\medskip\noindent\textbf{Welche Wiederholung nehmen wir?}
Wir wählen \(j\) als den \emph{ersten} Exponenten, dessen Potenz schon
vorher auftrat. Dann sind \(a^1,\ldots,a^{j-1}\) paarweise verschieden.
Darum gibt es genau einen Exponenten \(i<j\) mit \(a^i=a^j\).
Wir setzen
\[
  \mu=i,\qquad \lambda=j-i.
\]
Beide Zahlen sind positiv. \(\mu\) heißt \emph{Index}, \(\lambda\)
\emph{Periode}. Bis unmittelbar vor den Exponenten \(\mu+\lambda\)
sehen wir lauter verschiedene Werte; bei diesem Exponenten kehrt die
Folge zu \(a^\mu\) zurück. So erhält man das eindeutige
Index--Perioden-Paar aus
\FormulaRefAuto{FiniteSemigroupIndexPeriod}[thm].

Die Wahl des ersten Wiederholungsendpunkts ist wesentlich. Aus irgendeiner
Gleichheit \(a^i=a^j\) allein könnte man keine minimale Periode ablesen.
Im Beispiel gilt etwa auch \(a^2=a^6\), obwohl die tatsächliche Periode
\(2\) und nicht \(4\) ist. Die erste Wiederholung ist \(a^4=a^2\),
also \((\mu,\lambda)=(2,2)\).

\medskip\noindent\textbf{Warum bleibt die Wiederholung erhalten?}
Aus \(a^{\mu+\lambda}=a^\mu\) folgt für \(t\geq1\) durch
Multiplikation mit \(a^t\)
\[
  a^{\mu+\lambda+t}
  =a^{\mu+\lambda}a^t
  =a^\mu a^t
  =a^{\mu+t}.
\]
Für \(t=0\) ist dies bereits die Ausgangsgleichheit; hier wird kein
\(a^0\) benutzt. Insgesamt gilt somit
\[
  a^{r+\lambda}=a^r\qquad(r\geq\mu).
\]
Das ist \FormulaRefAuto{FiniteSemigroupPowerEventuallyPeriodic}[thm].
Eine gleiche Potenz wird mit demselben nächsten Faktor weitergeführt.
Die Assoziativität sorgt dafür, dass dieses Weiterführen mit dem
Additionsgesetz vereinbar ist.

Die Potenzfolge besteht daher aus einem Anlauf
\(a^1,\ldots,a^{\mu-1}\) und einem wiederkehrenden Kreis
\(a^\mu,\ldots,a^{\mu+\lambda-1}\). Beim Index \(1\) ist der Anlauf
leer. Jeder spätere Exponent lässt sich durch wiederholtes Abziehen
von \(\lambda\) auf den Kreis zurückführen. Weil die ersten
\(\mu+\lambda-1\) Werte verschieden sind, hat die von \(a\) erzeugte
Unterhalbgruppe genau
\[
  |\langle a\rangle|=\mu+\lambda-1
\]
Elemente; dies ist \FormulaRefAuto{FiniteSemigroupMonogenicSize}[thm].
Im Beispiel ist \(\langle a\rangle=S\), weil \(a,b,c\) als positive
Potenzen von \(a\) vorkommen. In einer beliebigen Halbgruppe kann
\(\langle a\rangle\) dagegen eine echte Teilmenge sein. Die Potenzfolge
eines Elements bestimmt dann nur die Produkte innerhalb dieser
Unterhalbgruppe, nicht die ganze Grundtafel.

\section*{Aus einem Kreis entsteht ein Idempotentes}

Ein Element \(e\) heißt \emph{idempotent}, wenn \(e^2=e\).
Jede Potenzfolge in einer endlichen Halbgruppe enthält ein solches
Element. Dafür brauchen wir keine neue Fallunterscheidung über die
Tafel: Wähle zu ihrem Index--Perioden-Paar einen positiven Exponenten
\(n\geq\mu\), der ein Vielfaches von \(\lambda\) ist. Zum Beispiel
funktioniert immer \(n=\mu\lambda\).

Ab \(n\) dürfen wir beliebig oft um eine Periode weiterschieben.
Weil auch \(n\) selbst ein Periodenvielfaches ist, folgt
\[
  (a^n)^2=a^{2n}=a^n.
\]
Dies beweist \FormulaRefAuto{FiniteSemigroupIdempotentPower}[thm].
Bei nichtleerem \(S\) wählen wir irgendein \(a\in S\) und erhalten
damit wenigstens ein idempotentes Element, wie in
\FormulaRefAuto{FiniteSemigroupIdempotentExists}[thm].

Im Dreierbeispiel dürfen wir \(n=2\) wählen, denn \(2\geq\mu\)
und \(2\) ist ein Vielfaches der Periode. Es ergibt sich
\(a^2=b\) und \(b^2=b\). Die stets mögliche Wahl
\(n=\mu\lambda=4\) liefert denselben Wert \(a^4=b\).
\(b\) ist auf \(\{b,c\}\) neutral, weil \(bb=b\) und \(bc=cb=c\).
Auf ganz \(S\) ist es nicht neutral, denn \(ba=c\ne a\).
Die Existenz eines Idempotenten allein macht aus einer Halbgruppe
also noch kein Monoid mit einem globalen neutralen Element.

\section*{Warum es ein kleinstes nichtleeres Ideal gibt}
\hypertarget{finite-semigroups.reading.ideals}{}

Ein nichtleeres \emph{zweiseitiges Ideal} ist eine nichtleere Teilmenge
\(I\subseteq S\), aus der man durch Multiplikation mit beliebigen
Elementen von \(S\) nicht herauskommt:
\[
  SI\subseteq I,\qquad IS\subseteq I.
\]
Hier bedeutet \(SI=\{xy\mid x\in S,y\in I\}\), entsprechend für
\(IS\). Das ist eine andere Frage als die nach den Potenzen eines
einzelnen Elements: Beide Seiten des Produkts dürfen sich ändern.

Ist \(S\ne\varnothing\), ist \(S\) selbst ein nichtleeres Ideal.
Zudem können zwei nichtleere Ideale \(I,J\) nicht disjunkt sein.
Wähle \(x\in I\) und \(y\in J\). Dann liegt \(xy\) wegen
\(IS\subseteq I\) in \(I\), und wegen \(SJ\subseteq J\) in \(J\).
Also ist \(I\cap J\ne\varnothing\). Ihr Durchschnitt ist wiederum
ein Ideal: Wer in beiden Mengen liegt, bleibt bei Multiplikation
in beiden Mengen.

Weil \(S\) endlich ist, gibt es nur endlich viele seiner Teilmengen
und somit nur endlich viele nichtleere Ideale. Wähle unter ihnen
ein bezüglich Inklusion minimales Ideal \(K\). Für jedes weitere
nichtleere Ideal \(J\) ist \(K\cap J\) ein nichtleeres Ideal in \(K\).
Minimalität erzwingt \(K\cap J=K\), also \(K\subseteq J\).
\(K\) ist daher sogar das \emph{kleinste} nichtleere Ideal:
Es liegt in jedem anderen. Zwei kleinste Ideale müssten ineinander
enthalten und damit gleich sein. Dies ist der Beweisgedanke von
\FormulaRefAuto{FiniteSemigroupLeastIdeal}[thm].

\begin{FiniteSemigroupsReadingExample}{Das kleinste Ideal im Dreierbeispiel}
In der Tafel stehen nur \(b\) und \(c\) als Produktwerte. Also ist
\(K=\{b,c\}\) ein nichtleeres Ideal. Jedes nichtleere Ideal enthält
mindestens eines der Elemente \(a,b,c\). Enthält es \(a\), enthält es
wegen \(a^2=b\) auch \(b\). Enthält es \(b\), enthält es wegen
\(ba=c\) auch \(c\). Enthält es \(c\), enthält es wegen \(ca=b\)
auch \(b\). In jedem Fall enthält es \(K\). Damit ist \(K\) das
kleinste nichtleere Ideal.

Hier fällt das Ideal mit dem Kreis der Potenzfolge von \(a\) zusammen.
Der allgemeine Idealbeweis benötigt diese besondere Darstellung nicht.
Er funktioniert auch in Halbgruppen, die kein einzelnes Element erzeugt.
\end{FiniteSemigroupsReadingExample}

\section*{Dieselbe Tafel unter anderen Namen}
\hypertarget{finite-semigroups.reading.tables}{}

Ein \emph{Isomorphismus} zwischen Halbgruppen \(S\) und \(T\) ist eine
Bijektion \(f\colon S\to T\) mit
\[
  f(xy)=f(x)f(y)\qquad(x,y\in S).
\]
Die Namen der Elemente dürfen wechseln; die Multiplikation muss dabei
mitwechseln. Insbesondere werden idempotente Elemente auf idempotente
Elemente abgebildet, denn \(f(e)^2=f(e^2)=f(e)\).
Induktiv gilt auch \(f(a^r)=f(a)^r\), sodass die ganze Potenzfolge
und ihr Index--Perioden-Paar erhalten bleiben.

Wir nummerieren einen Träger mit \(n\) Elementen durch
\(S_n=\{0,\ldots,n-1\}\). Eine Bijektion \(e\colon S_n\to S\)
transportiert die Multiplikation auf eine Zahlentafel
\[
  m(i,j)=e^{-1}\bigl(e(i)e(j)\bigr).
\]
Der Eintrag sagt, \emph{welche Nummer} das Produkt trägt.
Die Nummerierung \(e\) ist selbst ein Isomorphismus zwischen dieser
Tafel und der Ausgangshalbgruppe; Assoziativität und Endlichkeit
bleiben erhalten. Dieser Strukturtransport ist
\FormulaRefAuto{FiniteSemigroupMultiplicationTableTransport}[thm].

Für das Beispiel wählen wir \(e(0)=a,e(1)=b,e(2)=c\). Die Tafel lautet
\[
  \begin{array}{c|ccc}
    m&0&1&2\\\hline
    0&1&2&1\\
    1&2&1&2\\
    2&1&2&1
  \end{array}.
\]
Die Zahl \(0\) ist hier nur ein Name. Sie bezeichnet weder ein
Nullelement noch ein neutrales Element.

\medskip\noindent\textbf{Was muss beim Umbenennen geändert werden?}
Eine Permutation \(\pi\colon S_n\to S_n\) ist eine Umbenennung der
Nummern. Die neue Tafel \(d\) muss für alle \(i,j\) erfüllen
\[
  d\bigl(\pi(i),\pi(j)\bigr)=\pi\bigl(m(i,j)\bigr).
  \tag{U}\label{finite-semigroups-reading-relabel}
\]
Man ändert die Zeilenköpfe, die Spaltenköpfe \emph{und} jeden
Ergebnisnamen. Nur Zeilen und Spalten zu vertauschen würde im
Allgemeinen eine andere Operation erzeugen.

Vertauschen wir im Beispiel \(0\) und \(1\) und lassen \(2\) fest,
so erhalten wir
\[
  \begin{array}{c|ccc}
    d&0&1&2\\\hline
    0&0&2&2\\
    1&2&0&0\\
    2&2&0&0
  \end{array}.
\]
Zum Beispiel wird \(m(0,0)=1\) zur Gleichheit \(d(1,1)=0\).
Das idempotente Element trägt jetzt die Nummer \(0\). Alle neun
Einträge folgen aus derselben Regel~\eqref{finite-semigroups-reading-relabel}.

Diese Regel ist nicht nur notwendig. Erfüllt eine Bijektion \(\pi\)
die Gleichungen für alle Felder, dann erhält sie gerade jedes Produkt
und ist deshalb ein Isomorphismus. Damit ist
\FormulaRefAuto{FiniteSemigroupTableRelabelCriterion}[thm] bewiesen.
Zwei nummerierte Tafeln gleicher Größe sind also genau dann isomorph,
wenn eine gemeinsame Umbenennung aller Namen die eine in die andere
überführt. Die Menge aller umbenannten Tafeln heißt ihr \emph{Orbit}.
\FormulaRefAuto{FiniteSemigroupRelabelOrbitsAreIsomorphismClasses}[thm]
sagt, dass diese Orbits genau die Isomorphieklassen sind.

\section*{Ein vollständiger Code für gegebene Grundtafeln}

Eine nummerierte Tafel kann man zeilenweise als Tupel der Länge \(n^2\)
lesen. Bei \(m\) ist dies
\[
  w_3(m)=(1,2,1,\ 2,1,2,\ 1,2,1),
\]
bei \(d\) dagegen
\[
  w_3(d)=(0,2,2,\ 2,0,0,\ 2,0,0).
\]
Verschiedene Tupel können somit dieselbe Halbgruppe beschreiben.
Aber gleiche Tupel beschreiben dieselbe nummerierte Tafel: Jede
Position speichert genau ein festgelegtes Tafelfeld. Die Tupel sind
keine Dezimalzahlen; die Trennzeichen erhalten ihre Bedeutung auch
bei Nummern mit mehreren Ziffern.

Ein nummerierungsunabhängiger Code besteht aus der Trägergröße und
der Menge \emph{aller} Tafeltupel, die durch Umnummerierungen entstehen.
So ist \(\operatorname{can}(S)\) in
\FormulaRefAuto{FiniteSemigroupCanonicalTableCodeDef}[def] definiert.
Für die Berechnung kann man stattdessen das lexikographisch kleinste
Tupel dieses Orbits speichern: Man vergleicht die ersten verschiedenen
Einträge zweier Tupel und nimmt das mit dem kleineren Eintrag.
Die Trägergröße wird weiterhin mitgespeichert.

\medskip\noindent\textbf{Warum ist der Code vollständig?}
Ist \(f\colon S\to T\) ein Isomorphismus, führt jede Nummerierung
\(e\) von \(S\) durch \(f\circ e\) zu einer Nummerierung von \(T\)
mit genau denselben Tafeleinträgen. Mit \(f^{-1}\) erhält man die
Rückrichtung; die beiden Mengen von Tafeltupeln sind gleich.

Haben umgekehrt beide Halbgruppen denselben Code, dann haben sie
gleich viele Elemente. Ein Tafeltupel irgendeiner Nummerierung
\(e_S\) von \(S\) tritt auch bei einer Nummerierung \(e_T\) von
\(T\) auf. Die nummerierten Tafeln sind daher gleich. Die Bijektion
\[
  f=e_T\circ e_S^{-1}
\]
erhält jedes Produkt und ist der gesuchte Isomorphismus. Das sind
beide Richtungen von
\FormulaRefAuto{FiniteSemigroupCanonicalTableCodeComplete}[thm].

Die komprimierte Variante mit dem kleinsten Tupel funktioniert ebenso:
Haben zwei Orbits dasselbe Minimum, haben sie eine Tafel gemeinsam.
Jede weitere Tafel im einen Orbit entsteht aus dieser durch Umbenennung
und liegt deshalb auch im anderen Orbit. Die Orbits sind gleich.
Für unser Dreierbeispiel ist das kleinste Tupel
\[
  (0,1,1,\ 1,0,0,\ 1,0,0).
\]
Es entsteht, wenn \(b\) die Nummer \(0\), \(c\) die Nummer \(1\)
und \(a\) die Nummer \(2\) erhält.

Es gibt nur \(n!\) Permutationen, und zu jeder werden \(n^2\)
Einträge berechnet. Das Verfahren ist daher endlich. Vorher können
wir einfache Merkmale vergleichen, etwa Kommutativität oder die
Anzahl idempotenter Elemente. Ein Unterschied schließt Isomorphie
sofort aus. Gleiche Merkmale allein rechtfertigen sie noch nicht;
dafür dient der vollständige Code oder eine konkrete Bijektion mit
allen Produktgleichungen.

\section*{Wenn Teilmengen zu den neuen Elementen werden}
\hypertarget{finite-semigroups.reading.subsets}{}

Zu einer Halbgruppe \(S\) gehört ihre \emph{Potenzhalbgruppe}
\(\PowNE S\). Ihre Elemente sind die nichtleeren Teilmengen von
\(S\); ihre Multiplikation lautet
\[
  XY=\{xy\mid x\in X,\ y\in Y\}.
\]
Man bildet sämtliche Einzelprodukte und entfernt doppelte Werte.
Für nichtleere \(X,Y\) ist auch \(XY\) nichtleer, weil sich wenigstens
ein Paar von Faktoren wählen lässt. Die Grundassoziativität liefert
\((XY)Z=X(YZ)\): Beide Mengen bestehen aus genau den Werten
\((xy)z=x(yz)\) mit \(x\in X,y\in Y,z\in Z\).
Band~28 beweist diese Halbgruppenstruktur; ihre Endlichkeit steht in
\FormulaRefAuto{FinitePowerSemigroup}[thm].

Bei \(n\) Grundelementen gibt es \(2^n-1\) nichtleere Teilmengen.
Für jedes Grundelement entscheiden wir, ob es enthalten ist; von
den \(2^n\) Entscheidungen schließen wir nur die leere Menge aus.
Im Dreierbeispiel sind es sieben. Wir nummerieren sie wie folgt:
\[
  \begin{array}{c|ccccccc}
    p&0&1&2&3&4&5&6\\\hline
    \nu(p)&\{a\}&\{b\}&\{c\}&\{a,b\}&\{a,c\}&\{b,c\}&S
  \end{array}.
\]
Dies sind sieben neue \emph{Elemente der Potenzhalbgruppe}, auch wenn
wir sie weiterhin als Mengen schreiben.

\medskip\noindent\textbf{Wie wird ein Potenztafelfeld berechnet?}
Für \(X=\{a,b\}\) und \(Y=\{a,c\}\) lesen wir das Teilrechteck
mit den Zeilen \(a,b\) und den Spalten \(a,c\):
\[
  XY=\{aa,ac,ba,bc\}=\{b,b,c,c\}=\{b,c\}.
\]
Das Ergebnis trägt in der Potenztafel die Nummer \(5\), also
\(m_{\mathcal P}(3,4)=5\).
Eine Zweiermenge auf jeder Seite liefert zwar vier Einzelprodukte,
aber hier nur zwei verschiedene Werte. Im selben Eingabeblock
liefern \(\{a,c\}\{a,c\}=\{b\}\) sogar nur einen Wert. Die Größen
der Eingabemengen legen die Ergebnisgröße daher nicht fest.

Die vollständige Potenztafel des Beispiels ist
\[
  \begin{array}{c|ccc|ccc|c}
    m_{\mathcal P}&0&1&2&3&4&5&6\\\hline
    0&1&2&1&5&1&5&5\\
    1&2&1&2&5&2&5&5\\
    2&1&2&1&5&1&5&5\\\hline
    3&5&5&5&5&5&5&5\\
    4&1&2&1&5&1&5&5\\
    5&5&5&5&5&5&5&5\\\hline
    6&5&5&5&5&5&5&5
  \end{array}.
\]
Die Linien unterscheiden Einermengen, Zweiermengen und die Dreiermenge
in \emph{dieser gewählten Darstellung}. Im ersten
\(3\times3\)-Block steht die ursprüngliche Grundtafel, denn
\(\{x\}\{y\}=\{xy\}\). Alle übrigen Felder entstehen nach derselben
Teilrechteckregel. In einer so konstruierten und beschrifteten Tafel
kennen wir also die Einermengenkopie der Grundhalbgruppe ganz genau.

\section*{Warum die Einermengenkopie nicht automatisch erkennbar ist}
\hypertarget{finite-semigroups.reading.reconstruction}{}

Bei der Rekonstruktionsfrage ist nur ein \emph{abstrakter}
Isomorphismus
\[
  \Phi\colon\PowNE S\longrightarrow\PowNE T,\qquad
  \Phi(XY)=\Phi(X)\Phi(Y)
\]
gegeben. Er erhält die Multiplikation und ist bijektiv. Dass er
Teilmengengrößen, Vereinigungen oder Inklusionen erhält, gehört nicht
zu seinen Voraussetzungen. Die Linien und Mengenbeschriftungen der
vorigen Tafel sind somit keine zusätzliche algebraische Information,
die wir bei diesem Isomorphietest verlangen dürften.

Würde \(\Phi\) die Einermengen bijektiv auf Einermengen abbilden,
könnten wir \(\Phi(\{x\})=\{f(x)\}\) setzen. Dann wäre \(f\) bijektiv,
und aus
\[
  \{f(xy)\}=\Phi(\{xy\})
  =\Phi(\{x\})\Phi(\{y\})
  =\{f(x)f(y)\}
\]
folgte \(f(xy)=f(x)f(y)\). Genau diese bedingte Rekonstruktion
beweist \FormulaRefAuto{PowerSemigroupSingletonLayerReconstruction}[thm].
Die Bedingung selbst muss aber eigens gerechtfertigt werden.

\begin{FiniteSemigroupsReadingExample}{Schon unser Beispiel vertauscht eine Einermenge mit einer Zweiermenge}
In der Potenztafel sind Zeile \(0\) und Zeile \(4\) gleich, ebenso
Spalte \(0\) und Spalte \(4\). Die Mengen \(\{a\}\) und \(\{a,c\}\)
haben also dasselbe Multiplikationsverhalten gegenüber jeder
nichtleeren Testmenge, links wie rechts.

Vertausche diese beiden Mengen und lasse alle anderen Mengen fest.
Diese Abbildung \(\Phi\) ist bijektiv. Das Vertauschen eines Faktors
ändert wegen der gleichen Zeilen und Spalten kein Produkt. Jeder
Produktwert ist zudem eine der Mengen \(\{b\},\{c\},\{b,c\}\),
also wird jeder Produktwert von \(\Phi\) festgehalten. Damit gilt
für alle \(X,Y\)
\[
  \Phi(XY)=XY=\Phi(X)\Phi(Y).
\]
\(\Phi\) ist ein Automorphismus der Potenzhalbgruppe, obwohl es die
Einermenge \(\{a\}\) auf die Zweiermenge \(\{a,c\}\) abbildet.
\end{FiniteSemigroupsReadingExample}

Das Beispiel widerlegt die Behauptung, \emph{jeder gegebene}
Potenzhalbgruppenisomorphismus müsse Einermengen erhalten. Es liefert
kein Paar nichtisomorpher Grundhalbgruppen: Wir haben die
Potenzhalbgruppe von \(S\) mit sich selbst verglichen. Für eine
Rekonstruktion darf ein geeigneter Grundisomorphismus deshalb auch
auf anderem Weg entstehen; er muss sich nicht aus dem gerade
gegebenen \(\Phi\) auf Einermengen ablesen lassen.

Die Zahl der Grundelemente ist hingegen bereits gesichert. Auch wenn
bei einem Partner \(T\) zunächst keine Endlichkeit vorausgesetzt wird,
folgt sie aus dem Potenzisomorphismus: \(\PowNE T\) ist dann genauso
endlich wie \(\PowNE S\), und \(t\mapsto\{t\}\) bildet \(T\)
injektiv in diese endliche Menge ab. Für die nun endlichen \(S,T\) gilt
\[
  2^{|S|}-1=|\PowNE S|=|\PowNE T|=2^{|T|}-1.
\]
Da die Zweierpotenzen streng wachsen, folgt \(|S|=|T|\).
So können sieben Potenzhalbgruppenelemente nur von drei
Grundelementen stammen. Dies ist die Rückrechnung aus
\FormulaRefAuto{FiniteNonemptyPowerSetReflectsCardinality}[thm].
Sie bestimmt die Anzahl, ohne die einzelnen Einermengen zu bestimmen.

\section*{Rekonstruktion bei zwei Grundelementen}

Wie viel eine endliche vollständige Prüfung leisten kann, zeigt die
Ordnung \(2\). Auf \(\{0,1\}\) besitzt eine binäre Tafel vier Felder,
für jedes gibt es zwei mögliche Werte. Es gibt also \(2^4=16\)
Kandidaten. Bei jedem werden alle acht Tripel auf Assoziativität
geprüft. Nach dieser Prüfung und dem Zusammenfassen umbenannter
Tafeln bleiben genau fünf Isomorphietypen; die vollständigen Tafeln
und ausgeschlossenen Fälle stehen in
\FormulaRefAuto{FinitePowerSemigroupReconstructionOrderTwo}[thm].

Die Potenzhalbgruppe hat in jedem Fall drei Elemente. Dennoch lässt
ihre abstrakte Multiplikation die fünf Grundtypen unterscheiden:
\begin{center}
\small
\begin{tabularx}{\linewidth}{@{}>{\raggedright\arraybackslash}Xc
  >{\raggedright\arraybackslash}X@{}}
\toprule
Grundregel&Idempotente in \(\PowNE S\)&Weiteres Kennzeichen\\
\midrule
Alle Produkte gleich \(0\)&1&konstante Multiplikation\\
\(xy=x\)&3&\(XY=X\) für alle \(X,Y\)\\
\(xy=y\)&3&\(XY=Y\) für alle \(X,Y\)\\
\(xy=\max\{x,y\}\)&3&kommutativ\\
\(xy=x+y\pmod2\)&2&kommutativ\\
\bottomrule
\end{tabularx}
\end{center}
In der konstanten Halbgruppe ist nur \(\{0\}\) idempotent.
In der Zweiergruppe sind die Einermenge des neutralen Elements
und die volle Grundmenge idempotent; die andere Einermenge ist es
nicht. In den drei übrigen Fällen ist jede nichtleere Teilmenge
idempotent, wie die jeweiligen Produktregeln unmittelbar zeigen.
Die Idempotentenanzahl trennt also den konstanten Fall und den
Gruppenfall ab. Unter den drei verbleibenden Fällen ist allein der
Halbverband kommutativ.

Auch Linksnull- und Rechtsnullregel bleiben unterscheidbar.
Gäbe es einen Isomorphismus \(F\) von der ersten Potenzhalbgruppe
zur zweiten, ergäbe sich für verschiedene \(X,Y\)
\[
  F(X)=F(XY)=F(X)F(Y)=F(Y),
\]
was ihrer Bijektivität widerspricht. Damit gehören zwei
zweielementige Halbgruppen mit isomorphen Potenzhalbgruppen zum
selben der fünf Grundtypen und sind selbst isomorph. Die Aussage
beruht auf einer erschöpfenden Prüfung und anschließenden
algebraischen Unterscheidung, nicht auf der Kennzeichnung von
Einermengen.

\section*{Was eine endliche Suche für eine feste Ordnung beweist}
\hypertarget{finite-semigroups.reading.search}{}

Die Grundtafeln lassen sich nicht nur vergleichen, sondern auch
vollständig erzeugen. Fixiere \(N\geq1\). Jede Operation auf
\(S_N\) besitzt \(N^2\) Felder mit je \(N\) möglichen Werten, also
gibt es \(N^{N^2}\) Kandidaten. Man könnte sie alle ausfüllen und
jedes Tripel prüfen. Praktischer ist es, die Assoziativität schon
während des Ausfüllens zu benutzen.

Für jedes Tripel müssen
\[
  m\bigl(m(i,j),k\bigr)=m\bigl(i,m(j,k)\bigr)
  \tag{A}\label{finite-semigroups-reading-associativity}
\]
gelten. Sind beide Seiten bereits bekannt und verschieden, kann
die aktuelle Teiltafel nicht fortgesetzt werden. Sind die beiden
inneren Werte und ein äußerer Wert bekannt, ist der andere äußere
Wert erzwungen. Sind die inneren Werte bekannt und die beiden
äußeren Felder noch unbelegt, müssen diese Felder jedenfalls
denselben künftigen Wert erhalten. Nach solchen Folgerungen wählt
man für ein noch offenes Feld jeden noch möglichen Wert und
setzt die Suche in diesem Zweig fort. Diese Verzweigung mit
Rückkehr bei einem Widerspruch heißt \emph{Backtracking}.

Warum kann dadurch keine Halbgruppe verloren gehen? Betrachte eine
beliebige vollständige assoziative Tafel. Bei jeder freien Wahl
enthält die Verzweigung auch den Wert, der in dieser Tafel steht.
Jede erzwungene Gleichheit gilt in ihr wegen
\eqref{finite-semigroups-reading-associativity}; entlang dieses Zweigs
kann daher kein wirklicher Widerspruch entstehen. So erreicht die
Suche auch diese Tafel. Umgekehrt erfüllt eine vollständig
ausgefüllte Ausgabe alle geprüften Tripel und ist assoziativ.
Die Suche endet, weil nur endlich viele Felder mit endlich vielen
Werten belegt werden können. Der rechnerische Anhang von Band~37
unterscheidet diese Tafelsuche ausdrücklich von der bloßen
Fortsetzung einer einzelnen Potenzfolge.

Um nur einen Vertreter jeder Isomorphieklasse auszugeben, behalten
wir anschließend jeweils den kanonischen Tafelcode. Dazu dient
der zuvor bewiesene vollständige Isomorphietest. Nun kann die
Rekonstruktionsfrage für genau diese Ordnung untersucht werden:

\begin{enumerate}
  \item Erzeuge einen Vertreter jeder Halbgruppenisomorphieklasse
    auf \(S_N\).
  \item Konstruiere seine Potenzhalbgruppe durch die Teilrechteckregel.
  \item Berechne deren kanonischen Tafelcode als abstrakte Halbgruppe
    mit \(2^N-1\) Elementen.
  \item Vergleiche diese Potenztafelcodes zwischen den verschiedenen
    Grundvertretern.
\end{enumerate}

Taucht derselbe Potenztafelcode bei zwei verschiedenen
Grundvertretern auf, ist ein Gegenbeispiel dieser Ordnung gefunden:
Die Potenzhalbgruppen sind isomorph, die Grundhalbgruppen sind es
nicht. Taucht kein Code doppelt auf, ist die Rekonstruktion für
\emph{alle Halbgruppen der Ordnung \(N\)} bewiesen.
Die Rückrechnung der Grundelementzahl sichert, dass ein möglicher
Partner einer solchen endlichen Halbgruppe dieselbe Ordnung hat.

\medskip\noindent\textbf{Welche Umnummerierungen sind beim dritten Schritt erlaubt?}
Alle Permutationen der Potenzhalbgruppenelemente. Wir dürfen die
Einermengen oder die Teilmengengrößen nicht als festzuhaltende
Markierungen vorgeben. Unser Dreierbeispiel zeigt einen zulässigen
Automorphismus, der eine solche Markierung durchbricht.
Ein Test mit diesen Markierungen würde deshalb eine zusätzliche
Bedingung prüfen, die in der Rekonstruktionsfrage nicht gegeben ist.
Als Vorfilter sind dagegen Merkmale zulässig, die aus der
Multiplikation folgen, etwa Idempotenz. Jeder Isomorphismus muss
sie tatsächlich erhalten.

Die Endlichkeit ist hier eine Aussage über Terminierung, nicht
über eine günstige Laufzeit. Die Potenztafel besitzt bereits
\((2^N-1)^2\) Felder. Eine naive Kanonisierung würde bis zu
\((2^N-1)!\) Umnummerierungen durchlaufen. Verbesserte Suchregeln
können Arbeit sparen; sie müssen dabei begründen, dass aus jeder
möglichen Isomorphieklasse wenigstens ein Vertreter erhalten bleibt.
Auch eine vollständige Prüfung vieler kleiner Ordnungen liefert
für sich genommen keinen Beweis für beliebige \(N\).

\section*{Wie die Lesefassungen zusammengehören}

Die endliche Grundtafel wird durch ihren kanonischen Tafelcode
vollständig bestimmt. Die Potenzhalbgruppe lässt sich aus dieser
Tafel ebenfalls vollständig berechnen. Beim Rückweg fehlt jedoch
zunächst die algebraische Kennzeichnung der Grundelemente unter
allen nichtleeren Teilmengen. Weder die Linien einer Darstellung
noch das Aussehen der Mengenbeschriftungen leisten diese Kennzeichnung.

Die gesonderte \FiniteSemigroupsNullproductsReadingLink{} behandelt
einen strukturellen Weg durch diese Schwierigkeit: Für eine endliche
Halbgruppe mit Nullelement und genau zwei Produktwerten
\(S^2=\{0,c\}\), \(c\ne0\), werden aus Nulltests und endlichen
Anzahlen genügend Informationen gewonnen, um einen Grundisomorphismus
zu konstruieren. Dabei muss der gegebene Potenzisomorphismus die
Einermengen nicht selbst erkennen. Unser Dreierbeispiel besitzt
ebenfalls zwei Produktwerte, \(S^2=\{b,c\}\), aber kein Nullelement;
die Voraussetzung des Nullproduktsatzes ist daher nicht erfüllt.

Die vorliegende Fassung erklärt die allgemeinen endlichen Werkzeuge
und die Grenzen ihres unmittelbaren Einsatzes. Die Nullprodukt-Fassung
führt einen besonderen Rekonstruktionsbeweis zusammenhängend aus.
Die genaue Formulierung der allgemeinen Rekonstruktionsfrage und
ihre mit Quellen belegte Einordnung bleiben in Band~37. Alle hier
verlinkten nummerierten Aussagen führen dorthin oder zu den
benötigten Grundlagen in früheren Bänden.

\noindent\FiniteSemigroupsMainLink{}
\enspace\textperiodcentered\enspace
\FiniteSemigroupsNullproductsReadingLink{}
